\begin{longtable}{|p{2.9cm}|p{3.6cm}|p{3.6cm}|p{3.8cm}|}

\caption{Các đánh đổi trong quyết định thiết kế hệ thống
\label{tab:design-tradeoff}} \\

\hline
\textbf{Quyết định} &
\textbf{Thuộc tính thu được} &
\textbf{Thuộc tính phải hy sinh} &
\textbf{Cơ sở của lựa chọn} \\
\hline
\endfirsthead

\hline
\textbf{Quyết định} &
\textbf{Thuộc tính thu được} &
\textbf{Thuộc tính phải hy sinh} &
\textbf{Cơ sở của lựa chọn} \\
\hline
\endhead

\hline
\endfoot

\hline
\endlastfoot

Sổ đăng ký tin cậy trên chuỗi công khai thay cho cơ sở dữ liệu
&
Bên xác minh đọc được trạng thái thẩm quyền trực tiếp, không cần quan hệ trước với tổ chức nào; lịch sử quản trị không thể sửa
&
Chi phí giao dịch cho mỗi thao tác quản trị; trạng thái thẩm quyền trở thành thông tin công khai
&
Nếu sổ đăng ký là cơ sở dữ liệu thì bên xác minh lại phải tin tổ chức vận hành nó --- tức chỉ di chuyển sự phụ thuộc chứ không loại bỏ
\\ \hline

Một nguồn niềm tin gốc duy nhất
&
Mô hình quản trị đơn giản, trách nhiệm rõ ràng, phù hợp phạm vi kiểm chứng khái niệm
&
Tồn tại một điểm tập trung quyền lực quản trị
&
Phản ánh đúng cấu trúc thể chế hiện có của ĐHQGHN; kiến trúc phân lớp cho phép chuyển sang mô hình liên hợp mà không thay đổi ba lớp dưới
\\ \hline

Ví không lưu ký
&
Chủ thể là bên duy nhất ký được bản trình bày của mình, nhờ đó bước kiểm tra chữ ký của người trình bày có giá trị thật
&
Trải nghiệm phức tạp hơn; mất cụm từ khôi phục là mất quyền truy cập, không có cơ chế đặt lại
&
Mô hình lưu ký làm thay đổi tính chất an toàn mà hệ thống cung cấp, không chỉ làm giảm tính tự chủ về nguyên tắc
\\ \hline

Tách dữ liệu trên chuỗi và ngoài chuỗi
&
Chi phí lưu trữ thấp; không có dữ liệu cá nhân trên hạ tầng không thể xóa
&
Phụ thuộc vào tính khả dụng dài hạn của lớp lưu trữ ngoài chuỗi và cơ chế ghim
&
Ghi dữ liệu cá nhân lên hạ tầng bất biến khiến không thể thực hiện quyền của chủ thể dữ liệu theo quy định pháp luật
\\ \hline

Trạng thái theo từng bản ghi thay cho danh sách bit
&
Truy vấn trạng thái tức thời; thống nhất hạ tầng của ba sổ đăng ký trên cùng một lớp
&
Mỗi lần thu hồi là một giao dịch riêng; mỗi lần tra cứu để lại dấu vết về chứng thực được kiểm tra
&
Đơn giản hóa vận hành ở quy mô thử nghiệm; đây là điểm cần chuyển đổi khi mở rộng, và chỉ ảnh hưởng tới lớp chứng thực
\\ \hline

Chuẩn hóa JSON thay cho chuẩn hóa tập dữ liệu RDF
&
Triển khai được bằng mã ngắn, đồng nhất được giữa nhiều ngôn ngữ, dễ kiểm chứng tính tất định
&
Chỉ bảo đảm tương đương ở mức cú pháp; bên xác minh ngoài hệ thống cần mã kiểm tra riêng
&
Tương đương ngữ nghĩa của RDF yêu cầu bên xác minh triển khai toàn bộ phép xử lý JSON-LD, chi phí không tương xứng ở phạm vi hiện tại
\\ \hline

Giao diện dịch vụ theo kiểu REST thay cho giao thức truyền thông điệp giữa các tác nhân
&
Giảm đáng kể độ phức tạp triển khai và tích hợp
&
Khả năng tương tác giữa các hệ sinh thái thấp hơn kiến trúc tham chiếu đầy đủ
&
Phù hợp mục tiêu kiểm chứng mô hình quản trị; việc bổ sung giao thức truyền thông điệp không làm thay đổi kết luận về khung quản trị
\\ \hline

Giới hạn chuỗi ủy quyền ở hai cấp
&
Trách nhiệm luôn truy vết được tới một đơn vị cụ thể
&
Không biểu diễn được cấu trúc tổ chức sâu hơn hai cấp
&
Chuỗi ủy quyền sâu không giới hạn khiến việc truy vết trách nhiệm trở nên bất khả thi, trong khi lợi ích không rõ ràng
\\ \hline

\end{longtable}
